\section{\thunderbolt as a Payment Channel}\label{sec:protocol}

Having established the core cryptographic primitive in \autoref{sec:coreproof}, we now formalize \thunderbolt as a \emph{payment channel}---a long-lived off-chain object that is anchored by a fixed on-chain UTXO and evolves through a sequence of off-chain state updates. This section introduces the participants, the channel lifecycle, and the key state invariants. We deliberately use the same terminology as the classical payment-channel literature~\cite{poon2016lightning, decker2015fast} so that the differences are easy to see.

\subsection{Participants and Setup}\label{sec:participants}

A \thunderbolt channel has three participants roles:

\begin{itemize}
  \item The \textbf{Holder} (initially $\mathsf{H}_0$) is the party that currently owns the UTXO off-chain. Over the lifetime of the channel, the holder role is handed from one user to the next: $\mathsf{H}_0 \to \mathsf{H}_1 \to \mathsf{H}_2 \to \cdots$, where each $\mathsf{H}_s$ is an independent user.
  \item The \textbf{Committee} $\mathsf{C}$ is a fixed (but possibly threshold) party that maintains the complementary share $y_s$ of the lock key. $\mathsf{C}$ never owns the UTXO and never learns the holder secret $x_s$; its sole role is to preserve the fixed-key invariant~\eqref{eq:invariant} via the compensating ``subtract'' update.
  \item The \textbf{Chain} plays the role of a passive verifier: it sees exactly two transactions over the channel's lifetime (Setup and Finalize), regardless of how many holder transitions occur in between.
\end{itemize}

\paragraph{Setup (on-chain).} At channel open, the initial holder $\mathsf{H}_0$ and the Committee jointly derive the fixed public key $P = (x_0 + y_0) G$, where $x_0$ is $\mathsf{H}_0$'s initial holder secret and $y_0$ is the Committee's initial share. $\mathsf{H}_0$ funds a UTXO on the Bitcoin chain locked to a script whose spending condition enforces knowledge of the discrete log of $P$ (e.g., a standard BIP340 Taproot output~\cite{bip340}). This is the \emph{only} on-chain operation at open. The initial public anchor is $X_0 = x_0 H$, which is stored in the channel's off-chain state.

\paragraph{Off-chain state at step $s$.} The channel state at step $s \geq 0$ consists of:
\begin{itemize}
  \item the public anchor $X_s$,
  \item the holder ciphertext $T_s$ (a Paillier ciphertext under the current holder's public key, encrypting the holder secret $x_s$),
  \item the committee share $y_s$ (held by $\mathsf{C}$; in a threshold committee, held collectively via Shamir secret sharing),
  \item the chain of transfer proofs $\{(\pi^{(s')}_{\mathrm{inv}}, \pi^{(s')}_{\mathrm{xfer}})\}_{s' \leq s}$.
\end{itemize}

The fixed-key invariant $(x_s + y_s) G = P$ holds at every step by \autoref{thm:correct-step}.

\subsection{Channel Lifecycle}\label{sec:lifecycle}

A \thunderbolt channel passes through three phases: \textbf{Open}, \textbf{Transfer} (zero or more times), and \textbf{Finalize}.

\paragraph{Open.} $\mathsf{H}_0$ and $\mathsf{C}$ jointly generate $(x_0, y_0)$ and publish the on-chain UTXO under $P = (x_0 + y_0) G$. $\mathsf{H}_0$ initializes $T_0 = \Enc_{pk_{\mathsf{H}_0}}(x_0)$, $X_0 = x_0 H$.

\paragraph{Transfer.} To transfer ownership from $\mathsf{H}_s$ (Alice) to $\mathsf{H}_{s+1}$ (Bob):
\begin{enumerate}
  \item Bob chooses a fresh refresh value $r_s$, computes $D_s = r_s H$, $\mathrm{ct}_B = \Enc_{pk_B}(r_s)$, $\mathrm{ct}_C = \Enc_{pk_C}(r_s)$, and produces $\pi_{\mathrm{inv}}$ (Core Proof~1). Bob publishes the invoice tuple $(D_s, \mathrm{ct}_B, \mathrm{ct}_C, \pi_{\mathrm{inv}})$.
  \item Any verifier (including Alice, the Committee, and any future recipient) can check $\pi_{\mathrm{inv}}$ with respect to the public parameters. Alice need not contact Bob directly: she simply fetches the invoice.
  \item Alice computes $T_{s+1} = \Enc_{pk_B}(x_s) \oplus \mathrm{ct}_B$, the new anchor $X_{s+1} = X_s + D_s$, and produces $\pi_{\mathrm{xfer}}$ (Core Proof~2). She publishes the transfer tuple $(T_{s+1}, X_{s+1}, \pi_{\mathrm{xfer}})$.
  \item The Committee decrypts $\mathrm{ct}_C$ to recover $r_s$ and updates its share to $y_{s+1} = y_s - r_s$.
  \item Bob decrypts $T_{s+1}$ to recover $x_{s+1} = x_s + r_s$.
\end{enumerate}
No step in this phase requires Alice and Bob to be online simultaneously. No step touches the Bitcoin chain.

\paragraph{Finalize.} When the current holder $\mathsf{H}_s$ wishes to withdraw on-chain, she and $\mathsf{C}$ jointly produce a signature on the closing transaction under $P = (x_s + y_s)G$, using any standard two-party ECDSA or Schnorr protocol~\cite{lindell2017fast, komlo2021frost}. Because the on-chain key has not changed since Open, the Bitcoin node verifies the signature exactly as it would any other Taproot spend.

\subsection{Key Invariants}\label{sec:invariants}

The \thunderbolt protocol preserves the following invariants at every step $s$, assuming all transfer proofs verify:

\begin{enumerate}
  \item \textbf{Fixed-key invariant:} $(x_s + y_s) G = P$. The on-chain lock key never changes.
  \item \textbf{Anchor correctness:} $X_s = x_s H$. The public anchor is the discrete-log image of the current holder secret.
  \item \textbf{Ciphertext correctness:} $T_s = \Enc_{pk_{\mathsf{H}_s}}(x_s)$. The current holder ciphertext encrypts the current holder secret under the current holder's key.
  \item \textbf{Proof chain:} For every step $s' < s$, both $\pi^{(s')}_{\mathrm{inv}}$ and $\pi^{(s')}_{\mathrm{xfer}}$ verify.
\end{enumerate}

Invariants 1--3 follow from \autoref{thm:correct-step} by induction on $s$, and invariant 4 is a public-check by any verifier.

\subsection{Comparison with Lightning Channels}\label{sec:vs-lightning}

\thunderbolt and Lightning both realize the notion of an off-chain channel anchored by an on-chain UTXO, but the resulting usage models are strikingly different.

\begin{itemize}
  \item \textbf{Channel graph.} Lightning is a routed graph: every payment from $u$ to $v$ must traverse a path of pre-funded bidirectional channels with sufficient capacity in both directions. \thunderbolt has no graph: ownership moves atomically from one holder to one recipient, regardless of any prior connectivity.
  \item \textbf{Liquidity.} Lightning channels have directional liquidity; they must be rebalanced. \thunderbolt transfers move the entire UTXO and carry no notion of directional capacity.
  \item \textbf{Online assumptions.} Lightning requires both endpoints to be online at the moment of payment for HTLC setup. \thunderbolt transfers are asynchronous: Alice and Bob never need to be online at the same time. The Committee, however, must be available to decrypt $\mathrm{ct}_C$ and update its share; this is analogous to Lightning's ``watchtower'' requirement but weaker in that the Committee processes refresh updates, not fraud detection.
  \item \textbf{On-chain footprint per transfer sequence.} Both amortize on-chain cost over many off-chain transfers, but Lightning requires one on-chain open and close per \emph{channel} between each pair of endpoints, while \thunderbolt requires one open and close per \emph{UTXO} across an arbitrary chain of holders.
  \item \textbf{Trust model.} Lightning is trustless modulo watchtowers. \thunderbolt relies on the Committee to preserve the committee share; in the threshold-committee instantiation via VESS, this trust assumption is reduced to standard $(t, n)$-threshold honesty.
\end{itemize}

A full quantitative comparison appears in \autoref{sec:eval}.
